% Licence CC-BY-SA
% produit par ox223252
% document de base pour modification
\documentclass[a4paper]{report}
\usepackage{verbatim}   % permet d'utiliser verbatim
\usepackage{tcolorbox}  % permet d'afficher des boites autour de certain element
\usepackage[nobottomtitles*]{titlesec}   % permet de modifier les titres, l'option permet d'eviter le titre en fin de page
\usepackage{graphicx}   % permet d'inclure des images
						% for '\resizebox` macro
\usepackage[utf8]{inputenc}     % permet de changer l'encodage
\usepackage[T1]{fontenc}        % et d'utiliser le package frencais
\usepackage[french]{babel}     % utilisation des charactères français
\usepackage{float}      % permet d'utiliser \begin{figure}[H]
\usepackage{lscape}     % affichegage de page paysage
\usepackage{pmboxdraw}  % pour afficher les charactere de tree
\usepackage{amssymb,pdftexcmds,xparse}
\usepackage{cite}       % permet les citations
\usepackage[babel=true]{csquotes}
\usepackage[colorlinks=true]{hyperref}   % href link, pdf editables : https://la-bibliotex.fr/2020/04/11/creer-des-pdf-editables-avec-latex/
\usepackage{pst-uml}    % cf https://www.mathforu.com/hors-programme/latex-dessin-geometrique-en-latex-avec-pstricks/
\usepackage{ulem}       % barrer du text 
\usepackage{ccicons}    % icons : https://la-bibliotex.fr/2019/02/10/des-icones-a-infini-dans-latex/
\usepackage{subfig}     % sous figure 1a, 1b, ... : https://la-bibliotex.fr/2019/01/31/package-latex-decouverte-1/
\usepackage[final]{pdfpages}    % insertion de page de pdf : https://la-bibliotex.fr/2019/01/31/package-latex-decouverte-1/
\usepackage[yyyymmdd,hhmmss]{datetime} % ajout de la date \today \currenttime
\usepackage[acronym,nomain]{glossaries} % use glossaries-package
\usepackage{siunitx} % dessins de graphiques depuis banque de données
\usepackage{tikz}
\usepackage{pgfplots}
\usepackage{pdftricks} % convert your existing postscript files to PDF : alow to use tikz width pdflatex ?
\usepackage{ulem} % barrer du text \sout{Texte à barrer} \xout{Texte à hachurer}
\usepackage{array, multirow}
\usepackage{fancyhdr} % entête
\usepackage{fontawesome5} % font awesome

\usepackage[european,straightvoltages,siunitx]{circuitikz} % circuit elecs
\usetikzlibrary{babel}

\usepackage[left=3.5cm,right=3.5cm,top=2cm,bottom=2cm]{geometry} % changement des marges

% pip install pygments
\usepackage{minted}     % code en couleur
\setlength\partopsep{-\topsep}
\addtolength\partopsep{-\parskip}
% \usemintedstyle{rrt}

\titleformat{\chapter}[display]{\bfseries\Huge}{}{0em}{}[\vspace{-0.5em}]
\def\changemargin#1#2{\list{}{\rightmargin#2\leftmargin#1}\item[]}
\let\endchangemargin=\endlist

\usepackage{cmbright}

%
% permet de faire des checbox
%
\tikzset{
	box/.style={
		minimum size=0.225cm,
		inner sep=0pt,
		draw,
	},
	insert mark/.style={
		append after command={
			node[inner sep=0pt,#1]
			at (\tikzlastnode.center){$\checkmark$}
		}
	},
	insert bad mark/.style={
		append after command={
			[shorten <=\pgflinewidth,shorten >=\pgflinewidth]
			(\tikzlastnode.north west)edge[#1](\tikzlastnode.south east)
			(\tikzlastnode.south west)edge[#1](\tikzlastnode.north east)
		}
	},
}

\makeatletter \NewDocumentCommand{\tikzcheckmark}{O{} m}{
	\ifnum\pdf@strcmp{#2}{mark}=\z@
		\tikz[baseline=-0.5ex]\node[box,insert mark={#1},#1]{};
	\fi
	\ifnum\pdf@strcmp{#2}{bad mark}=\z@
		\tikz[baseline=-0.5ex]\node[box,insert bad mark={#1},#1]{};
	\fi
	\ifnum\pdf@strcmp{#2}{no mark}=\z@
		\tikz[baseline=-0.5ex]\node[box,#1]{};
	\fi
}
\makeatother

% \tikzcheckmark{no mark}
% \tikzcheckmark{mark}
% \tikzcheckmark{bad mark}

% permet de dessiner un code avec coloration syntaxique et encadrement
\newenvironment{cbox}[1]
{
	\def\temp{#1}\ifx\temp\empty
		\def\a{}
	\else
		\def\a{\caption{#1}}
	\fi

	\begin{figure}[H]
	\begin{changemargin}{-2.5cm}{-2.5cm}
	\begin{tcolorbox}
}
{
	\end{tcolorbox}
	\end{changemargin}
	\a
	\end{figure}
}

%
% environement pour include des mails
%
\newenvironment{mail}[4]
{
	\subsection{#1}
		\begin{changemargin}{-2.5cm}{-2.5cm}
		\begin{tcolorbox}
		\begin{itemize}
			\item From: {#2},
			\item To: {#3},
			\item Date: {#4},
		\end{itemize}
}
{
		\end{tcolorbox}
		\end{changemargin}
}

% permet d'afficher un extrait de document

\newsavebox{\auteurbm}

%
% environement pour citation à droite
%
% param : nom de l'auteur
\newenvironment{extraitR}[1]{\small\slshape
\savebox{\auteurbm}{\upshape\sffamily#1}
\begin{flushright}}{\\[4pt]\usebox{\auteurbm}
\end{flushright}\normalsize\upshape}

%
% environement pour citation à gauche
%
% param : nom de l'auteur
\newenvironment{extraitL}[1]{\small\slshape
\savebox{\auteurbm}{\upshape\sffamily#1}
\begin{flushleft}}{\\[4pt]\usebox{\auteurbm}
\end{flushleft}\normalsize\upshape}

\titleformat{\chapter}[display]{\bfseries\Huge}{}{0em}{}[\vspace{-0.5em}]
\def\changemargin#1#2{\list{}{\rightmargin#2\leftmargin#1}\item[]}
\let\endchangemargin=\endlist


%
% flow chart
%
\usetikzlibrary{shapes.geometric, arrows}

\tikzstyle{startstop} = [rectangle, rounded corners, minimum width=3cm, minimum height=1cm,text centered, draw=black, fill=red!30]
\tikzstyle{io} = [trapezium, trapezium left angle=70, trapezium right angle=110, minimum width=3cm, minimum height=1cm, text centered, draw=black, fill=blue!30]
\tikzstyle{process} = [rectangle, minimum width=3cm, minimum height=1cm, text centered, text width=3cm, draw=black, fill=orange!30]
\tikzstyle{decision} = [diamond, minimum width=3cm, minimum height=1cm, text centered, draw=black, fill=green!30]
\tikzstyle{arrow} = [thick,->,>=stealth]

